Time-series foundation models have emerged as a powerful paradigm for zero-shot and few-shot temporal forecasting, with architectures such as Chronos~\cite{FastAccurateZeroshot2024} achieving remarkable generalisation across diverse domains.
A key enabler of this scalability is \emph{patch-based tokenisation}: rather than feeding individual scalar observations to a transformer, contiguous windows of samples are grouped into fixed-length patches, dramatically shortening the effective sequence length and reducing the quadratic cost of self-attention.
Despite its computational appeal, this design choice introduces a subtle but consequential inductive bias that has received little formal scrutiny.

The central concern of this work is what we term \emph{structural aliasing}: a representational failure that arises not from undersampling in the classical Nyquist–Shannon sense, but from the geometric interaction between the periodicity of an input signal and the fixed stride of the patch grid.
When the stride $S$ is commensurate with the signal period $T_0$, consecutive patches become identical or near-identical, collapsing the token sequence to a constant or low-diversity representation.
The model then suffers a loss of temporal observability that is invisible to standard sampling-theoretic guarantees — the signal is fully representable, yet the architecture cannot perceive it.

Beyond the signal-theoretic dimension, structural aliasing carries semantic consequences.
In real-world deployment scenarios, time-series signals encode domain concepts such as operational states, anomalies, control conditions, and recurring events.
A model that is blind to certain frequencies may therefore misclassify or entirely miss semantically meaningful patterns.
To make this connection explicit, we introduce an ontological layer that maps signal-level observability failures to domain-level concept distortions, and we evaluate the phenomenon through both behavioural metrics and internal representational probing.

This report is structured as follows.
Section~\ref{sec:one} formalises patch-based tokenisation as a linear operator and derives the conditions under which token identity — and hence observability loss — occurs.
Section~\ref{sec:two} characterises the resulting blind-spot frequencies as harmonics of the patch frequency $f_{\mathrm{patch}} = f_s / S$, and analyses the role of patch size, stride, and overlap.
Section~\ref{sec:point3} introduces the semantic layer through a formal ontology of the chosen application domain.
Section~\ref{sec:point4} enumerates the falsifiable hypotheses that drive our experimental programme.
Sections~\ref{sec:point5} and~\ref{sec:point6} describe the controlled synthetic experiments and the evaluation methodology at both the behavioural and representational levels.
Section~\ref{sec:point7} presents the Bayesian statistical analysis used to quantify uncertainty and compare competing hypotheses.
Section~\ref{sec:point8} investigates a principled mitigation strategy, and Section~\ref{sec:point9} situates the findings within an agentic AI setting, discussing how an autonomous agent could monitor, explain, and act upon structural aliasing risks.
